\documentclass[12pt]{article}
\usepackage[letterpaper, margin=1in]{geometry}
\usepackage{titlesec}
\usepackage{titling}
\usepackage{graphicx}
\usepackage{float}
\usepackage{url}
\usepackage{color}
\usepackage{siunitx}
\usepackage{natbib}
\usepackage{subcaption}
\usepackage{amsmath,amssymb,amsfonts}
\usepackage[colorlinks=true, urlcolor=blue, linkcolor=blue, citecolor=blue]{hyperref}

% Title formatting courtesy of Damein Koon
\pretitle{\LARGE\bfseries}
\posttitle{\par\vspace{2em}}
\preauthor{\large}
\postauthor{\par\vspace{1em}}
\predate{\large}
\postdate{\par\vspace{3em}}

% Corrected title hooks for vertical and horizontal centering
\renewcommand\maketitlehooka{\null\vfill\begin{center}}
\renewcommand\maketitlehookd{\end{center}\vfill\null}

\title{Machine Learning Regression on Synthetic Gravitational-Wave Inspirals}
\author{Krupa Pothiwala \\[0.5em]
February 2026}
\date{}

\begin{document}
\maketitle
\begin{abstract}
This project develops a machine learning surrogate model for gravitational-wave binary inspirals. A synthetic dataset of 1,000 samples was generated using the equations of orbital mechanics, the quadrupole approximation, and chirp mass theory. A Random Forest regression model was trained to recover chirp mass and merger time from strain time-series features. Feature importance analysis revealed that the model learned algebraic shortcuts rather than physical relationships, highlighting the fundamental challenges of gravitational-wave parameter estimation from raw strain data.
\end{abstract}
\thispagestyle{empty}
\newpage

\section{Background Information}
\label{info}
The goal of this project is to develop a machine learning surrogate model for gravitational-wave inspirals by generating a synthetic dataset of 1,000 samples using the equations of orbital mechanics, quadrupole approximation, and chirp mass theory. Each sample consists of a 17-dimensional strain time-series, initial separation $r_0$, distance to observer $d$, and reduced mass $\mu$ as features paired with target labels for merger time, and chirp mass. Classical regression models will be trained to learn the mapping from strain features to binary parameters, creating a computationally efficient learned approximation of the underlying physics.

\subsection{Quadrupole Approximation}
A quadrupole moment describes the shape of a mass distribution, particularly how elongated or asymmetric it is. As two masses orbit each other, this shape rotates and evolves over time, and it is this changing configuration that radiates gravitational waves.
\\An easier way to visualize this is to imagine how the shape of a dumbbell changes as it rotates or evolves. 
\\
\textbf{Why is it important?}
Gravitational waves are generated by a time-varying quadrupole moment: a second-order spatial moment of the mass distribution \cite{CreightonGW}. When a binary system (two masses) orbits, the quadrupole moment changes periodically, and this changing quadrupole moment is what radiates gravitational waves.\\
\\
\textbf{Why use the quadrupole approximation?}
The quadrupole approximation relies on several simplifying assumptions about the system: 
These assumptions are:
\begin{itemize}
    \item The system is in the weak-field regime (gravity is not too strong)
    \item Velocities are non-relativistic (much less than the speed of light)
    \item The distance scale over which gravity changes is much larger than the size of the system
\end{itemize}


\section{Data Generation}
\label{generation}
Instead of directly using the quadrupole moment formula, I found an expression for the amplitude of the gravitational waves in \cite{CreightonGW}. I used this to calculate the amplitude of the gravitational waves.
\subsection{Physical Parameters}
\begin{enumerate}
    \item $m_1$ and $m_2$: Mass of the two objects. In the range of  $1.2 - 50$  $M_\odot$. Where $M_\odot$ is the solar mass \SI{1.989e30}{kg}. This limit includes neutron stars and most of the stellar-mass black holes. However, it does not include supermassive black holes.
    \item $r_0$: Separation radius which is the initial distance between the two objects. In the range of \SI{5e6}{km} to \SI{1e8}{km}. 
    \item $d$: Distance to the observer. In the range of $10$kpc to $10^5$kpc. Later converted to meters to keep everything in SI.
    \item $G=$ \SI{6.674e-11}{m^3/kg\,s^2} is the gravitational constant.
    \item $c=$ \SI{3e8}{m/s} is the speed of light.
\end{enumerate}

\subsection{Orbital Decay}
Orbital decay shows how the orbits of the two masses evolve with time.
\\The orbital decay equation from \cite{Peters1964} (Eq.~5.6) gives:
\begin{align}
    \frac{dr}{dt} = -\frac{64}{5} \frac{G^3 m_1 m_2 (m_1 + m_2)}{c^5 r^3}
    \label{eq:Decay}
\end{align}
This expression assumes a circular binary orbit with negligible eccentricity ($e=0$), which is typical for compact binaries whose eccentricity has been radiated away by gravitational-wave emission.

\subsection{Orbital Evolution}
The Ordinary Differential Equation was solved numerically using the \texttt{solve\_ivp} routine from the \texttt{scipy.integrate} package with the \texttt{RK45} method.

\subsection{Orbital Frequency}
Orbital frequency shows how fast the two masses are orbiting each other, and below I will derive it from first principles:

\begin{align}
    F_g &= F_{cen} \implies \frac{Gm_1m_2}{r^2} = m_1 \omega^2 r_1 \\
    r_1 &= \frac{m_2}{m_1+m_2}r 
    \implies 
    \frac{Gm_1m_2}{r^2} = \omega^2 \frac{m_1m_2}{m_1+m_2}r \\
    \omega^2 &= \frac{G(m_1 + m_2)}{r^3}
    \implies f_{orb} = \frac{1}{2\pi} 
    \sqrt{\frac{G(m_1 + m_2)}{r^3}} \label{eq:frequency}
\end{align}


\subsection{Strain Amplitude}
Amplitude of the gravitational waves also known as strain $h$. The strain is a dimensionless quantity: the fractional change in length. When a gravitational wave passes Earth, it stretches/compresses distances by a tiny fraction $h$.\\
\\The orbital (and gravitational-wave) frequency sweeps up in a "chirp" characterized by the chirp mass $\mathcal{M}$ \cite{CreightonGW}. The chirp mass also determines the frequency evolution of the gravitational wave signal emitted during a binary's inspiral phase \cite{ChirpMassWiki}. Basically, it is the parameter that controls how fast the binary spirals inward.
    \begin{align}
        \mathcal{M} \equiv \frac{(m_1m_2)^{3/5}}{(m_1+m_2)^{1/5}}
        \label{eq:Chirp}
    \end{align}
$\omega_r$ is the orbital angular frequency.
    \begin{align}
        \omega_r \equiv 2\pi f_{orb}
        \label{eq:angular}
    \end{align}
The equation below is taken from \cite{CreightonGW}. This expression represents the leading-order scaling of the gravitational-wave strain amplitude in the quadrupole approximation.
    \begin{align}
        h(t) &= \frac{4(G\mathcal{M})^{5/3}\omega_r^{2/3}}{c^4d}
        \label{eq:strain}
    \end{align}

\subsection{Time to merger}
The merger time measures how long it takes for two orbiting masses to completely spiral in and collide from a given starting separation $r_0$.
\\The equation for the merger time is derived from equation \eqref{eq:Decay}. Below I will show the derivation.
To simplify the equation \eqref{eq:Decay} I will define, $M = m_1m_2(m_1+m_2)$. \\
\begin{align}
     r^{3} dr &= \frac{-64}{5} \frac{G^3M}{c^5} dt\\
     \int_{r_0}^{0} r^{3} dr &= \int_0^{t_m}\frac{-64}{5} \frac{G^3M}{c^5} dt \\
     \frac{-r_0^4}{4} &= \frac{-64}{5} \frac{G^3M}{c^5} t_m \\
     t_m &=\frac{5c^5r_0^4}{256G^3M}
\end{align}
Here $t_m$ is time to merger. Let's substitute $M$ and get the final answer.
\begin{align}
    t_m &= \frac{5}{256}\frac{c^5r_0^4}{G^3m_1m_2(m_1+m_2)}
    \label{eq:merger}
\end{align}

\section{Data Preparation}
\label{dataprep}
Since the dataset was synthetically generated, no values were missing. Due to the values of my data being huge in SI units the values were transformed using a logarithmic scale to reduce the magnitude of the values. However, the values were converted back to linear space before calculating the mean absolute error. \\
The dataset was split into training and test sets using a 75–25 ratio. To prevent the data from being randomly shuffled on each run, I added a fixed random seed.\\
\\The project required the data set to have 20 features/dimensions and 1000 data points.\\
\\The features included: 
\begin{itemize}
    \item noisy strain values at 17 time stamps: $h_{noisy}(t=1),\dots,h_{noisy}(t=17)$,
    \item distance to the observer $d$, 
    \item initial separation between masses $r_0$,
    \item reduced mass: $\mu = \frac{m_1m_2}{m_1+m_2}$.
\end{itemize}
The 17 strain values and the rest of the values created the 20 features as per the requirements of the project.
\\
\\The targets included:
\begin{itemize}
    \item chirp Mass $\mathcal{M}$ \eqref{eq:Chirp}, 
    \item time to merger $t_m$ \eqref{eq:merger}\\
\end{itemize}

\section{Regression Model}
\label{model}
A Random Forest is an ensemble method, meaning that it combines many individual models (decision trees) to produce a better overall prediction than any single tree \cite{IBMRandomForest}.\\
\\My reasons for choosing this method:
\begin{itemize}
    \item The relationship between strain features and chirp mass/merger time is highly nonlinear: Random Forests handle this naturally without requiring polynomial feature engineering.
    \item My feature set mixes very different physical quantities (frequencies, distances, strain values) on vastly different scales: Random Forests are largely scale-invariant.
    \item With only 1,000 samples, overfitting is a real risk; ensemble averaging mitigates this.
\end{itemize}

\section{Results}
\label{sec:results}
The performance of the Random Forest regression model was evaluated using mean absolute error (MAE) and the coefficient of determination ($R^2$) on both the training and test datasets.\\
\\The training set had MAE of $4.1374 \times 10^{29}$ and the $R^2$ value of $0.8703$. \\
\\The test set had the MAE of $1.0793 \times 10^{30}$ and the $R^2$ value $0.5904$\\
\\While the MAE appears large, this reflects the scale of the physical quantities involved. The $R^2$ values of $0.87$ (train) and $0.59$ (test) suggest that the model is memorizing the $r_0$ values on training data but struggles to interpolate accurately on unseen $r_0$ values, especially at the extremes of the range as seen in Figure 1.
\begin{figure}[H]
    \centering
    \includegraphics[width=17cm]{predicted_vs_actual.png}
    \caption{Predicted vs. Actual values for both targets}
    \label{fig:}
\end{figure}
Looking at the results it seems that the model is learning. So, I removed the physical parameters that were used in the derivation of the strain just to see if the model would worsen. So instead of 17 strain values I increased it to 20 values. See Figure 2 for results.\\
\begin{figure}[H]
    \centering
    \includegraphics[width=17cm]{Only_h/predicted_vs_actual.png}
    \caption{Only strain: Predicted vs. Actual values for both targets}
    \label{fig:}
\end{figure}
As shown in Figure 2, the model does not appear to learn any meaningful patterns from the data. To investigate this further, I examined which parameters the model relied on most heavily in both cases.\\
\\The \texttt{scikit-learn} implementation of \texttt{RandomForestRegression} has an inbuilt feature called feature importance that lets the user see which feature the model relies on most when making predictions. It’s essentially a way to measure how much each feature contributes to reducing prediction error in the forest.
\begin{figure}[H]
\centering
    \begin{subfigure}{0.49\textwidth}
        \includegraphics[width=\textwidth]{Only_h/feature_importance.png}
        \caption{Strain-only features}
    \end{subfigure}
    \hfill
    \begin{subfigure}{0.49\textwidth}
        \includegraphics[width=\textwidth]{feature_importance.png}
        \caption{Full feature set}
    \end{subfigure}
\caption{Random Forest feature importances for chirp mass and time to merger prediction}
\label{fig:importance}
\end{figure}
Note that the near-zero and apparently negative values in Figure 3a are a plotting artefact of axis scaling; all importances are effectively zero, consistent with the model having found no learnable signal.\\
\\Examination of the feature importances revealed that the model was not utilizing the strain values in its predictions. Instead, it relied primarily on parameters that are directly related to the target quantity and appear in the analytical derivation, effectively reconstructing an algebraic relationship between these variables. The reduced mass enters implicitly through the definition of the chirp mass in equation \eqref{eq:Chirp}, while $r_0$ appears explicitly in equation \eqref{eq:merger}. Consequently, the model exploits these analytic relationships rather than learning patterns from the strain data, suggesting that the strain values alone contain insufficient accessible signal for this model to recover the target parameters.

\section{Conclusion}
\label{conclusion}
This project investigated whether a Random Forest regression model could recover gravitational wave binary parameters from synthetic inspiral data. Two different experiments were conducted with progressively restricted feature sets, and each revealed something important about how the model was actually learning.\\
\\In the first experiment, including derived quantities such as $r_0$, $d$, and $\mu$ as features produced decent predictions ($R^2 \approx 0.87$ for training and $R^2 \approx 0.59$ for test). However, as discussed in Section~\ref{sec:results}, feature importance  analysis revealed that this performance was not the result of genuine physical learning. Rather than extracting information from the strain signal, the model exploited the fact that $r_0$ and $\mu$ appear in the merger time (Equation 
\eqref{eq:merger}) and in chirp mass (Equation \eqref{eq:Chirp}), respectively, effectively reconstructing an algebraic relationship between the input features and the target quantities.\\
\\In the second experiment, restricting features to only the noisy strain time-series caused the model to fail entirely. The flat horizontal predictions in Figure 2 are the classic signature of a model defaulting to predicting the mean value, having found no learnable signal.\\
\\This progression demonstrates why real gravitational wave parameter estimation is genuinely difficult. Facilities like LIGO do not recover chirp mass and merger time by feeding raw strain into a simple regressor. 
Instead, they use matched filtering, which is cross-correlating the detector signal against thousands of precomputed waveform templates, then combined with Bayesian inference frameworks such as \texttt{LALInference}~\cite{Veitch2015} or \texttt{Bilby}~\cite{Ashton2019} to produce posterior probability distributions over parameters. A Random Forest operating on 17 or 20 noisy time samples cannot replicate this.\\
\\Next step could be to generate a much denser strain time-series (hundreds of time steps spanning the full inspiral), normalize the strain by distance to remove $d$ as an implicit shortcut, and explore other regression methods that could extract meaningful parameter information from the waveform morphology alone.

\newpage
\bibliographystyle{apalike}
\bibliography{references}
\end{document}